\section{The exact trace-norm obstruction in the PW exhaustion}
\label{sec:pw-trace-norm-convergence}

Strong convergence of the actual finite PW projections was proved
in MA-GP-001 and MA-CGP-003.  It does not by itself imply trace-norm
convergence of an uncut smoothed projection.  The following result
identifies that additional condition exactly, and proves its
sufficiency on the original weighted Hilbert space.

Use any of the actual rational or character families of MA-FT-001,
with \(0<\alpha<1/2\), \(s\ge0\), and
\[
 H=L^2(\mathbb R,w(x)\,dx),\qquad
 w(x)=e^{2\alpha|x|}(1+x^2)^{s/2}.
 \tag{MA-TN.1}
\]
Retain \(J_c,P_c,Q_c,P_c^Z\) from (MA-FT.3), and put
\[
 J=\overline{T_{1/2}\mathcal I}^{\,H},\qquad
 P_G=P_J,\qquad Z=J^\perp,\qquad P^Z=1-P_G.
 \tag{MA-TN.2}
\]
The full-strip lengths \(r_\rho^{\rm st}\), half-line lengths
\(r_\rho^+\), and nonnegative difference
\(\delta_\rho^-=r_\rho^+-r_\rho^{\rm st}\) are exactly (MA-FT.2).
In particular every length here is a length of an actual
finite primary quotient, with the boundary-jet thresholds
already proved in MA-WP and MA-GW.

Let \(U_tf(x)=f(x-t)\) on \(H\).  Write
\[
 \omega=\alpha+s/4,\qquad \nu>\omega,\qquad
 a_0=\nu-\omega,\qquad
 \ell(x)=(\log w)'(x)\quad(x\ne0).
 \tag{MA-TN.3}
\]
Thus \(|\ell|\le2\omega\).  On \(Z_c=P_c^ZH\) let \(B_c\)
be the adjoint quotient generator of MA-NZ, and on \(Z\)
let \(B\) be the adjoint of the actual full quotient
translation generator.  Extend their resolvents by zero:
\[
 R_c=(\nu-B_c)^{-1}P_c^Z,\qquad
 R=(\nu-B)^{-1}P^Z,\qquad
 A_c=\operatorname{Re}R_c,\quad A=\operatorname{Re}R.
 \tag{MA-TN.4}
\]
These are actual bounded operators on the single space \(H\).

\paragraph{MA-TN-001 (an exact positive trace defect).}
As \(c\to-\infty\), both \(R_c\to R\) and \(R_c^*\to R^*\)
strongly.  The operators \(A_c,A\) are positive and trace class,
and
\[
 \begin{aligned}
 \operatorname{Tr}A_c
   &=\sum_\rho r_\rho^+
                 \operatorname{Re}(\nu-\lambda_\rho)^{-1},\\
 \operatorname{Tr}A
   &=\sum_\rho r_\rho^{\rm st}
                 \operatorname{Re}(\nu-\lambda_\rho)^{-1},\\
 \lim_{c\to-\infty}\|A_c-A\|_1
   &=D_\nu:=
       \sum_\rho\delta_\rho^-
                 \operatorname{Re}(\nu-\lambda_\rho)^{-1},
       \qquad \lambda_\rho=\rho-\tfrac12.
 \end{aligned}
 \tag{MA-TN.5}
\]
All three sums converge absolutely.  In particular \(D_\nu=0\)
if and only if every \(\delta_\rho^-\) is zero.  If this happens,
\[
              \|R_c-R\|_2\longrightarrow0,\qquad
              \|R_c^2-R^2\|_1\longrightarrow0.
 \tag{MA-TN.6}
\]
There is no assertion that this arithmetic condition holds
without an additional argument.

\begin{proof}
\emph{Actual strong limits.}
The dense increasing PW exhaustion of MA-GP-001 or MA-CGP-003
gives \(Q_c\to P_G\) strongly, while \(P_c\to1\) strongly.
Consequently \(P_c^Z=P_c-Q_c\to P^Z\) strongly.  On the
supported half-line the adjoint of positive translation is
\(P_cU_t^*P_c\).  Since the actual \(J_c\) is invariant under
positive translation, its orthogonal complement in \(H_c\)
is invariant under this adjoint.  The adjoint semigroup,
extended by zero, is therefore \(P_cU_t^*P_c^Z\).
The full \(J\) is translation invariant, so the corresponding
full-space expression is \(U_t^*P^Z\).  The physical bound
\(\|U_t\|\le e^{\omega t}\), \(t\ge0\), gives the integrable
majorant in the vectorwise identities
\[
 R_c=\int_0^\infty e^{-\nu t}P_cU_t^*P_c^Z\,dt,\qquad
 R=\int_0^\infty e^{-\nu t}U_t^*P^Z\,dt .
 \tag{MA-TN.7}
\]
Strong convergence follows by dominated convergence for
each vector.  Taking adjoints in these integrals gives
integrands \(P_c^ZU_tP_c\) and \(P^ZU_t\), so the same
argument proves adjoint strong convergence.  All resolvent
norms are at most \(a_0^{-1}\).

\emph{Positive traces from the actual complete jet spaces.}
MA-NZ.6--.8 gives
\[
 a_0R_c^*R_c\le A_c,\qquad
 a_0R_cR_c^*\le A_c
 \tag{MA-TN.8}
\]
and the first trace in (MA-TN.5).  The identical energy
argument applies to \(R\).  For completeness, MA-WP.4 says
that \(Z\) is the closed span of the continuous retained
jet representatives
\[
       w(x)^{-1}x^j e^{\overline{\lambda_\rho}x},
                  \qquad 0\le j<r_\rho^{\rm st}.
 \tag{MA-TN.9}
\]
They are in the graph domain of the full adjoint generator,
whose expression is \(w^{-1}(w\,\cdot)'\).
Finite height spans are invariant and have eigenvalues
\(\overline{\lambda_\rho}\) with these exact lengths.
Their orthogonal projections increase strongly to \(P^Z\).
The trace of the compression of \(A\) to such a finite span
is the corresponding sum of
\(r_\rho^{\rm st}\operatorname{Re}(\nu-\lambda_\rho)^{-1}\).
Positivity and an orthonormal exhaustion, exactly as in
MA-NZ.8, prove that \(A\) is trace class and give its
asserted trace.

For every retained point \(\operatorname{Re}\lambda_\rho\le\alpha\).
Thus the summands are positive, and the real part of each
resolvent factor is \(O((1+|\operatorname{Im}\rho|^2)^{-1})\).
The fixed-function multiplicity-weighted local zero count
proves absolute convergence.  It also proves the assertion
about vanishing of \(D_\nu\), without assuming that any
discarded point exists.

\emph{The positive-operator trace-norm lemma.}
We prove the precise form needed here.  Suppose \(T_j,T\ge0\)
are trace class, \(T_j\to T\) strongly, and
\(\operatorname{Tr}T_j=S\) is constant, with
\(\operatorname{Tr}T\le S\).  Then
\[
             \|T_j-T\|_1\longrightarrow S-\operatorname{Tr}T.
 \tag{MA-TN.10}
\]
Choose a finite-rank orthogonal projection \(F\), and put
\(F^\perp=1-F\).  Strong convergence on the finite-dimensional
range of \(F\), and self-adjointness, show that the three
blocks with at least one \(F\) in the decomposition of
\(T_j-T\) tend to zero in trace norm.  For example
\((T_j-T)F\) has uniformly bounded rank and its
Hilbert--Schmidt norm tends to zero by summing over an
orthonormal basis of \(FH\).
The remaining block has the upper bound
\[
 \|F^\perp(T_j-T)F^\perp\|_1
 \le \operatorname{Tr}T_j-\operatorname{Tr}(FT_jF)
       +\operatorname{Tr}T-\operatorname{Tr}(FTF).
 \tag{MA-TN.11}
\]
First let \(j\) tend to the limit.  Then let \(F\) range
through a finite-rank exhaustion for the positive trace-class
operator \(T\).  The upper bound tends to
\(S-\operatorname{Tr}T\).  The reverse bound is
\(\|T_j-T\|_1\ge|\operatorname{Tr}(T_j-T)|\).
This proves (MA-TN.10); the proof applies to the parameter
\(c\to-\infty\), without discretizing its values.
Apply it to \(A_c,A\) to prove the last line of (MA-TN.5).

\emph{Hilbert--Schmidt convergence when the defect vanishes.}
Suppose \(D_\nu=0\), so \(A_c\to A\) in trace norm.
For a finite-rank orthogonal \(F\), (MA-TN.8) gives
\[
 \|R_c(1-F)\|_2^2
       \le a_0^{-1}\operatorname{Tr}((1-F)A_c).
 \tag{MA-TN.12}
\]
The analogous estimate holds for \(R\).
Choose \(F\) so that the right side for \(A\) is small;
trace-norm convergence makes it uniformly small for all
sufficiently negative \(c\).  On \(FH\), strong convergence
implies \(\|(R_c-R)F\|_2\to0\).  Splitting the right input
by \(F+(1-F)\) proves \(\|R_c-R\|_2\to0\).
The two Hilbert--Schmidt norms are bounded, so
\[
 \|R_c^2-R^2\|_1
   \le \|R_c\|_2\|R_c-R\|_2
             +\|R_c-R\|_2\|R\|_2\longrightarrow0 .
 \tag{MA-TN.13}
\]
This proves (MA-TN.6).
\end{proof}

\paragraph{MA-TN-002 (the actual ambient convergence criterion).}
The following conditions are equivalent:
\begin{enumerate}
 \item \(\delta_\rho^-=0\) for every \(\rho\);
 \item for every \(g\in C_c^\infty(\mathbb R)\),
 \(\|C_gP_c^Z-C_gP^Z\|_1\to0\);
 \item the convergence in condition 2 holds for every
 \(g\in C_c^\infty((0,\infty))\).
\end{enumerate}
These are actual ambient operators, not only quotient
compressions or equalities of scalar traces.
Strong convergence \(C_gP_c^Z\to C_gP^Z\) holds for every
fixed compact smooth \(g\) without condition 1.
If condition 1 fails, a causal test can be chosen for which
\[
 \liminf_{c\to-\infty}
       \|C_gP_c^Z-C_gP^Z\|_1
 \ge\left|\sum_\rho\delta_\rho^-G(\lambda_\rho)\right|>0,
 \qquad G(z)=\int g(t)e^{zt}\,dt .
 \tag{MA-TN.14}
\]
The test may also satisfy \(G(1/2)=G(-1/2)=0\).

\begin{proof}
Only sufficiency requires a new endpoint argument.
All operators whose norms are asserted to converge are
trace class by MA-AT-001 and MA-WP-003.
Assume condition 1.  Set \(k(x)=\nu-\ell(x)/2\);
it is a fixed bounded positive multiplier with \(k\ge a_0\).
For \(v\) in the first graph domain of the full half-line
adjoint generator, integration by parts gives
\[
 \operatorname{Re}\langle v,(\nu-\mathcal L_c)v\rangle_H
       =\int_c^\infty k(x)|v(x)|^2w(x)\,dx
                     +\tfrac12 w(c)|v(c+)|^2 .
 \tag{MA-TN.15}
\]
There is no cusp contribution in this first-order identity:
\(w\) is locally absolutely continuous.  The boundary
term at infinity vanishes.  Indeed
\(\sqrt w\,v\) belongs to ordinary \(H^1(c,\infty)\),
because \(\ell\) is bounded, and an \(H^1\cap L^2\)
function tends to zero at infinity.
The full-line version has neither endpoint term, by
the same argument at both ends.

Define the bounded scalar functional
\[
       e_c(f)=\sqrt{w(c)}(R_cf)(c+).
 \tag{MA-TN.16}
\]
Its boundedness follows from (MA-TN.15), applied to
\(v=R_cf\).  Polarization gives the exact operator identities
\[
 A_c=R_c^*M_kR_c+\tfrac12e_c^*e_c,\qquad
 A=R^*M_kR .
 \tag{MA-TN.17}
\]
MA-TN-001 gives convergence of the left sides in trace norm
and of \(R_c\) in Hilbert--Schmidt norm.  Consequently
\(R_c^*M_kR_c\to R^*M_kR\) in trace norm.  Subtracting
in (MA-TN.17) gives
\[
                         \|e_c\|^2\longrightarrow0 .
 \tag{MA-TN.18}
\]
This controls a moving boundary functional in its actual
norm, rather than by a nonuniform equivalence of weights.

Use the second-order bulk operator of MA-NZ.13:
\[
 \mathcal B_g=C_{g''}-2C_{g'}M_{\nu-\ell}
             +C_gM_{(\nu-\ell)^2-\ell'_{\rm reg}},
 \quad
 \ell'_{\rm reg}(x)=s(1-x^2)/(1+x^2)^2.
 \tag{MA-TN.19}
\]
It is bounded on \(H\) for every fixed compact smooth \(g\).
The exact second-order integration by parts in MA-NZ.14--.15,
now with \(r=g\), gives for \(c<0\)
\[
 C_gP_c^Z=\mathcal B_gR_c^2+\mathcal E_{c,g}R_c^2
                              -4\alpha\,g\otimes\kappa_c,
 \qquad
 \kappa_c(f)=(wR_c^2f)(0).
 \tag{MA-TN.20}
\]
The same calculation on the full line gives
\[
             C_gP^Z=\mathcal B_gR^2-4\alpha\,g\otimes\kappa,
 \qquad \kappa(f)=(wR^2f)(0).
 \tag{MA-TN.21}
\]
All graph domains and the cusp contribution are exactly
those checked in MA-NZ; in particular the jump is
\(\ell(0+)-\ell(0-)=4\alpha\).

Here is the moving endpoint in a form controlled by
(MA-TN.18).  Put \(g_c(x)=g(x-c)\).  Since
\(\mathcal L_cR_c^2=\nu R_c^2-R_c\), the endpoint
formula MA-NZ.14 becomes
\[
 \mathcal E_{c,g}R_c^2
  =\frac{(\nu-\ell(c))g_c-g_c'}{\sqrt{w(c)}}
                                  \otimes(e_cR_c)
       +\frac{g_c}{\sqrt{w(c)}}\otimes e_c .
 \tag{MA-TN.22}
\]
The sign and the coefficient follow by replacing
\((wR_c^2f)'(c)\) by
\(\sqrt{w(c)}(\nu e_cR_cf-e_cf)\) in that formula.
The norms of the two target vectors in (MA-TN.22)
are bounded uniformly in \(c<0\): on the fixed translated
support, \(w(c+t)/w(c)\le e^{2\omega|t|}\).
Together with \(\|R_c\|\le a_0^{-1}\), (MA-TN.18) proves
\(\|\mathcal E_{c,g}R_c^2\|_1\to0\).

The cusp functional has a fixed-vector description.
The actual half-line differential resolvent, squared, is
\[
 (\nu-\mathcal L_c)^{-2}h(x)
   =w(x)^{-1}\int_x^\infty
                  (y-x)e^{-\nu(y-x)}w(y)h(y)\,dy ,
       \qquad x\ge c .
 \tag{MA-TN.23}
\]
This follows either by two vectorwise semigroup integrals
or by composing their absolutely convergent kernels.
For \(c<0\), evaluation at \(x=0\) therefore says
\(\kappa_c(f)=\langle P_c^Zf,k_0\rangle_H\), where
\(k_0(y)=\mathbf1_{y\ge0}y e^{-\nu y}\), with the
inner product linear in its first argument.
The same formula holds with \(P^Z\) for \(\kappa\).
The vector \(k_0\) belongs to \(H\), since \(\nu>\alpha\).
Strong convergence of the orthogonal projections gives
\(\|\kappa_c-\kappa\|\to0\).

Now (MA-TN.6) controls the bulk difference in
(MA-TN.20)--(MA-TN.21), (MA-TN.22) controls the endpoint,
and the fixed-vector formula controls the cusp.
Their sum converges in trace norm.  This proves
condition 1 implies condition 2 for every compact smooth
test, without a condition at zero.

Condition 2 implies condition 3 by restriction.
For a causal test, the actual finite quotient trace
of MA-GW and the actual full-strip trace of MA-WP give
\[
 \operatorname{Tr}(C_gP_c^Z-C_gP^Z)
                     =\sum_\rho\delta_\rho^-G(\lambda_\rho).
 \tag{MA-TN.24}
\]
It is independent of \(c\).  The multiplicity-weighted
local count makes the series absolutely convergent.
If condition 3 holds, this sum vanishes for every
compact causal test.  MA-JP-004, applied to the
nonnegative integer divisor with coefficients
\(\delta_\rho^-\), then makes that divisor empty.
The completed nontrivial divisor in these primitive sectors
has no point at \(\rho=0\) or \(1\), as in MA-FT-002.
Thus any nonempty deficiency contains a point away from
the removed exponents \(\lambda=\pm1/2\).
The version of MA-JP-004 with those two polar moments
therefore gives the final assertion and (MA-TN.14).
Strong convergence, independently of
all these conditions, follows from \(P_c^Z\to P^Z\)
and boundedness of \(C_g\).
\end{proof}

\paragraph{Actual cutoff consequence and arithmetic scope.}
Under the equivalent conditions of MA-TN-002, let
\(M_n=M_{\chi_n}\) be any uniformly bounded compact
position multipliers converging strongly to one,
and let \(c_n\to-\infty\).  There is no required relative
rate.  Then the renormalized actual operators satisfy
\[
 \begin{aligned}
 M_nC_gQ_{c_n}-M_nC_gP_{c_n}
      &=-M_nC_gP_{c_n}^Z
                  \longrightarrow-C_gP^Z
                              \quad\hbox{in }\mathcal S_1(H).
 \end{aligned}
 \tag{MA-TN.25}
\]
Indeed split the difference into
\(-M_n(C_gP_{c_n}^Z-C_gP^Z)\) and
\(-(M_n-1)C_gP^Z\).  The first tends to zero by
MA-TN-002; the second does so by finite-rank approximation
of the fixed trace-class operator.  When \(g(0)=0\),
the subtracted compact-input term has trace zero
by MA-AT.12.  It is still an operator term and is not
discarded in the asserted trace-norm identity.
This does not assert trace-norm convergence of the
unrenormalized global operator, nor does it change the
different local geometric limits on different joint paths.
In particular the fast-path counterexample MA-JP-003
is unaffected.

The own-sector functional equation used in MA-FT-002
pairs the upper and lower deficiencies, including their
boundary lengths.  Thus the condition in MA-TN-002 is
equivalent to the exact upper-deficiency condition of
MA-FT-002.  For all positive widths it is the same sector's
RH statement.  The stipulated \(\mathsf A\) supplies it
for the rational and Dirichlet sectors explicitly as
zero-location input; \(M\) suffices for the specified
non-induced cubic Hecke sector.  The equivalence and
the actual conditional operator convergence above use
neither assertion in their proofs.

\paragraph{Trace-framework placement and what has changed.}
The coefficients \(\delta_\rho^-\) remain differences of
finite lengths in the spectral primary modules of MT-M18
and MA-WP.  Their additive trace contribution is the
finite-length \(K_0\) interpretation already used in
MA-EG and MA-FT, and the analytic traces belong to the
nuclear/Fredholm framework of Grothendieck, cote~1,
scans 79--80 and 111--112, as checked in
ALPOGE\_BRIDGES BR-TR-01/02.
The ordinary positive operator \(A_c\) is the real part
of an actual dissipative resolvent.  It is not a
faithful arithmetic star-action on the full jet module,
and it yields no weight, purity, or motivic positivity
theorem.  The Artin-motive interpretation of the character
sector remains that of MA-CH and MA-CQ, with SGA~1,
Expos\'e~V, \S7 and Deligne (1979), \S6.1--6.3 as the
already recorded exact placement.
This result identifies an operator topology consequence
of the existing arithmetic gate; it is not counted as a
strict reduction of that gate.
